\documentclass[11pt,a4paper]{article}
\usepackage{amsmath,amssymb,amsthm}
\usepackage[margin=1in]{geometry}
\usepackage{graphicx}
\usepackage{booktabs}
\usepackage[colorlinks=true,linkcolor=blue,citecolor=blue,urlcolor=blue]{hyperref}

\newtheorem{theorem}{Theorem}[section]
\newtheorem{proposition}[theorem]{Proposition}
\newtheorem{lemma}[theorem]{Lemma}
\newtheorem{corollary}[theorem]{Corollary}
\theoremstyle{definition}
\newtheorem{definition}[theorem]{Definition}
\newtheorem{hypothesis}{Hypothesis}
\renewcommand{\thehypothesis}{H\arabic{hypothesis}}
\theoremstyle{remark}
\newtheorem{remark}[theorem]{Remark}

\newcommand{\R}{\mathbb{R}}
\newcommand{\Z}{\mathbb{Z}}
\newcommand{\T}{\mathbb{T}}
\newcommand{\E}{\mathbb{E}}
\newcommand{\g}{\mathfrak{g}}
\newcommand{\Tr}{\operatorname{Tr}}
\newcommand{\tr}{\operatorname{tr}}
\newcommand{\Ad}{\operatorname{Ad}}
\newcommand{\dd}{\mathrm{d}}
\newcommand{\dstar}{\mathrm{d}^{*}}
\newcommand{\CalA}{\mathcal{A}}
\newcommand{\CalG}{\mathcal{G}}

\title{A Conditional Framework for the Four-Dimensional Euclidean Yang--Mills Measure on a Compact Polish Space of Wilson Loops via Strongly Parabolic Covariant Mollification}
\author{Max van Massey}
\date{\today}

\begin{document}
\maketitle

\begin{abstract}
We present a \emph{conditional} framework for the continuum limit of four-dimensional Euclidean Yang--Mills theory with compact matrix gauge group $G\le U(N)$. Gauge-invariant Wilson-loop data are placed in a compact Polish space $W_L$ built over a countable dense family of Lipschitz loops of bounded length and curvature, with the uniform topology. The weakly parabolic Yang--Mills flow is replaced by the strongly parabolic Yang--Mills--DeTurck flow $\partial_t B=-d_B^*F_B-d_Bd^*B$, whose drift gauge leaves every closed-loop trace invariant by cyclicity. We locate the critical Besov obstruction at regularity $-2-2\varepsilon$, derive the BPHZ constant $C_t=\tfrac12(K_{2t}-\delta_0)\otimes I_{\g}$, show that the mass-type counterterm vanishes under a Ward condition, and report a $16^4$ lattice-FFT measurement ($m_{\rm sim}=-0.092783$, $b_{\rm sim}=0.350591$). Non-triviality is encoded by a scale-response function along the two-loop trajectory and by two-sided exponential bounds on the fourth connected cumulant. Existence of the limit is proved only modulo three isolated analytic hypotheses \textup{(H1)--(H3)}.
\end{abstract}

\tableofcontents

%==========================================================================
\section{Introduction and Notation Framework}\label{sec:intro}
%==========================================================================

\subsection{Epistemological status}

The existence of a nontrivial four-dimensional Yang--Mills quantum field theory with a mass gap is open. This paper does \emph{not} claim a solution. It isolates a chain of reductions whose global conclusion holds \emph{conditionally} on three precisely stated, presently unproven analytic hypotheses:

\begin{hypothesis}[Uniform stochastic bounds]\label{h:H1}
Let $\Psi^{(\epsilon)}$ denote the regularised stochastic convolution of the DeTurck flow \eqref{eq:dtflow-stoch} at cutoff $\epsilon$, and let $\mathcal X^{(\epsilon)}$ be the finite collection of its BPHZ-renormalised Wick-type products. Then $\sup_{\epsilon}\E\|\mathcal X^{(\epsilon)}\|_{\mathcal C^{\alpha}}^{p}<\infty$ for all $p<\infty$, with the regularity indices $\alpha$ prescribed by Section~\ref{sec:besov}.
\end{hypothesis}

\begin{hypothesis}[Tightness in $W_L$]\label{h:H2}
The family of lattice Wilson-loop laws, restricted to the countable family $\mathcal D$ of Definition~\ref{def:WL}, is tight in $W_L$; equivalently the loop functions $\gamma\mapsto \tr\,\mathrm{Hol}_\gamma$ are, uniformly in the cutoff and with probability tending to one, equicontinuous with a fixed modulus $\omega_{L,K}$.
\end{hypothesis}

\begin{hypothesis}[Ward condition]\label{h:H3}
The renormalised DeTurck flow obeys the Ward condition $\mathrm W$ of Definition~\ref{def:ward}: the counterterms generated by the cutoff are covariant under $x$-dependent gauge transformations up to the DeTurck compensation, uniformly in the cutoff.
\end{hypothesis}

Every theorem below is of the form ``$\mathrm{(H1)}\wedge\mathrm{(H2)}\wedge\mathrm{(H3)}\Rightarrow\cdots$'' unless it is explicitly marked unconditional. The numerical material of Section~\ref{sec:bphz} is \emph{evidence} and is not part of any proof.

\subsection{Lie group structures}

Fix $N\ge 2$. Following the formalisation conventions of the accompanying Lean~4 development (namespace \texttt{YMDT}), $\mathrm{Mat}_N=\mathrm{Mat}_{N\times N}(\mathbb C)$ and a \emph{compact matrix Lie group} is a closed subgroup $G\le U(N)$ (structure \texttt{CompactMatrixLieGroup}: a \texttt{carrier} subgroup of \texttt{unitaryGroup} together with \texttt{isClosed}). Its Lie algebra is realised as
\begin{equation}
\g=\{X\in \mathrm{Mat}_N:\ \exists\, c:\R\to\mathrm{Mat}_N,\ c(0)=1,\ c(s)\in G\ \forall s,\ c'(0)=X\},
\end{equation}
which is a real subspace of the anti-Hermitian matrices with bracket $[X,Y]=XY-YX$. The invariant inner product is $\langle X,Y\rangle=-\Tr(XY)$, and $\{T_a\}_{a=1}^{\dim G}$ is an orthonormal basis, $[T_a,T_b]=f_{ab}{}^{c}T_c$, with $f_{abc}$ totally antisymmetric. The quadratic Casimir in a representation $\rho$ is $C_2(\rho)\,\mathbf 1=-\sum_a\rho(T_a)^2$.

\subsection{Torus, connections, gauge group}

Let $\T^4=(\R/\Z)^4$ with the periodic lift of the formalisation: $f:\R^4\to V$ is periodic if $f(x+e_\nu)=f(x)$ for each $\nu$. Smooth connections are $\CalA=\{A=A_\mu\,\dd x^\mu:\ A_\mu\in C^\infty(\T^4;\g)\}$ (\texttt{IsConn}), and the gauge group is $\CalG=C^\infty(\T^4;G)$ (\texttt{IsGauge}), acting by
\begin{equation}
A^g_\mu=g^{-1}A_\mu g+g^{-1}\partial_\mu g\qquad(g^{-1}=g^{\dagger}\ \text{for }G\le U(N)).
\end{equation}
The curvature and covariant derivative are, componentwise,
\begin{align}
F_{\mu\nu}&=\partial_\mu A_\nu-\partial_\nu A_\mu+[A_\mu,A_\nu],\\
D_\mu\varphi&=\partial_\mu\varphi+[A_\mu,\varphi],\\
(d_A^*F)_\nu&=-\sum_{\mu=1}^4 D_\mu F_{\mu\nu},\qquad d^*B=-\sum_{\mu=1}^4\partial_\mu B_\mu .
\end{align}
Gauge covariance: $F_{\mu\nu}(A^g)=g^{-1}F_{\mu\nu}(A)g$ and $(d_{A^g}^*F_{A^g})_\nu=g^{-1}(d_A^*F_A)_\nu\, g$.

\subsection{Wilson action and two-loop trajectory}

On a lattice $a\Z^4/L\Z^4$ with link variables $U_{x,\mu}\in G$ and plaquettes $U_p$, the Wilson action is
\begin{equation}
S_\beta(U)=\beta\sum_{p}\Bigl(1-\tfrac1N\mathrm{Re}\,\Tr\,U_p\Bigr),\qquad \beta=\frac{2N}{g_0^2}\quad(G=SU(N)),
\end{equation}
and the lattice measure is $\dd\mu_{a}(U)=Z^{-1}e^{-S_\beta(U)}\prod\dd U_{x,\mu}$ with Haar measure $\dd U$. Asymptotic freedom is encoded in the beta function $\mu\,\dd g/\dd\mu=-b_0g^3-b_1g^5+O(g^7)$,
\begin{equation}
b_0=\frac{11N}{3(4\pi)^2},\qquad b_1=\frac{34N^2}{3(4\pi)^4},
\end{equation}
whose integration gives the two-loop running
\begin{equation}\label{eq:2loop}
\frac1{g^2(\mu)}=2b_0\ln\frac{\mu}{\Lambda}+\frac{b_1}{b_0}\ln\ln\frac{\mu}{\Lambda}+O\!\Bigl(\frac{\ln\ln}{\ln}\Bigr).
\end{equation}
A continuum limit is a sequence $a\downarrow0$, $g_0=g_0(a)$ following \eqref{eq:2loop} at $\mu=1/a$. Note $\dd g^2/\dd\ln R=2b_0g^4+O(g^6)>0$ for $R=1/\mu$; this sign is used in Section~\ref{sec:nogo}.

\subsection{The compact Polish space $W_L$}

\begin{definition}[Loop class]\label{def:loops}
For $L,K>0$, let
\begin{equation}
\mathcal L_{L,K}=\bigl\{\gamma\in C^{1,1}([0,1];\T^4):\ \gamma(0)=\gamma(1),\ \|\gamma'\|_\infty\le L,\ \mathrm{Lip}(\gamma')\le K\bigr\},
\end{equation}
topologised by the $C^1$-norm $\|\gamma\|_{C^1}=\|\gamma\|_\infty+\|\gamma'\|_\infty$. (Length $\le L$ and curvature controlled by $K$.)
\end{definition}

\begin{lemma}\label{lem:compactloops}
$\mathcal L_{L,K}$ is a compact metrisable space.
\end{lemma}
\begin{proof}
Metrisability is inherited from the Banach space $C^1([0,1];\R^4)$ (the lift to $\R^4$ is fixed by $\gamma(0)\in[0,1)^4$ and the winding class, of which only finitely many are compatible with $\|\gamma'\|_\infty\le L$). Let $(\gamma_n)\subset\mathcal L_{L,K}$. The sequence $(\gamma_n')$ is uniformly bounded by $L$ and $K$-Lipschitz, hence equicontinuous; $(\gamma_n)$ is bounded and $L$-Lipschitz. By the Arzel\`a--Ascoli theorem there is a subsequence with $\gamma_{n_k}\to\gamma$ and $\gamma_{n_k}'\to\eta$ uniformly. Passing to the limit in $\gamma_{n_k}(s)=\gamma_{n_k}(0)+\int_0^s\gamma_{n_k}'$ gives $\gamma'=\eta$, so $\gamma\in C^1$ with $\|\gamma'\|_\infty\le L$ and $|\gamma'(s)-\gamma'(s')|\le K|s-s'|$; the closed condition passes to the limit. Hence the class is sequentially, thus (metric) compact.
\end{proof}

Since $\mathcal L_{L,K}$ is a compact metric space it is separable; fix a countable dense subset $\mathcal D\subset\mathcal L_{L,K}$.

\begin{definition}[Generalised connection space]\label{def:WL}
Let $\omega=\omega_{L,K}:[0,\infty)\to[0,\infty)$ be a fixed modulus of continuity. Put
\begin{equation}
W_L=\Bigl\{w:\mathcal D\to \overline{B_N(0)}\subset\mathbb C:\ |w(\gamma)-w(\gamma')|\le\omega(\|\gamma-\gamma'\|_{C^1})\ \ \forall \gamma,\gamma'\in\mathcal D\Bigr\},
\end{equation}
with the uniform topology $d_\infty(w,w')=\sup_{\gamma\in\mathcal D}|w(\gamma)-w'(\gamma)|$. The intended point of $W_L$ associated to a smooth connection $A$ is $w_A(\gamma)=\Tr\,\mathrm{Hol}_\gamma(A)$ (so $|w|\le N$).
\end{definition}

\begin{proposition}[$W_L$ is a compact Polish space]\label{prop:polish}
$(W_L,d_\infty)$ is compact and metrisable, hence Polish.
\end{proposition}
\begin{proof}
Each $w\in W_L$ extends uniquely to a function on the compact space $\mathcal L_{L,K}$ (the closure of $\mathcal D$), uniformly continuous with modulus $\omega$. Thus $W_L$ is identified with a subset of $C(\mathcal L_{L,K};\overline{B_N})$ which is uniformly bounded and equicontinuous; it is also closed in $d_\infty$ since both defining inequalities are preserved by uniform limits. By Arzel\`a--Ascoli it is compact. Subsets of the separable Banach space $C(\mathcal L_{L,K};\mathbb C)$ are metrisable and separable, so $W_L$ is a compact metric, hence complete separable metric, space.
\end{proof}

\begin{remark}[Why this is not the Tychonoff product]
The product $\prod_{\gamma\in\mathcal L_{L,K}}\overline{B_N}$ over the \emph{uncountable} loop family is compact by Tychonoff's theorem but is not metrisable (an uncountable product of nontrivial Hausdorff spaces fails first countability), hence not Polish, and its product topology is far weaker than the uniform topology; Prokhorov's theorem and regular conditional probabilities are unavailable there. Two ingredients restore Polishness: (i) indexing by the \emph{countable} family $\mathcal D$, so that the ambient product $\prod_{\mathcal D}\overline{B_N}$ is a countable product of compact metric spaces and therefore compact metrisable; (ii) the equicontinuity constraint, which makes $W_L$ a closed subset on which the uniform and product topologies coincide (on an equicontinuous set the pointwise topology on a dense set equals the uniform topology). The price is that equicontinuity with a fixed $\omega$ is a genuine restriction on the support, whose uniform validity for lattice measures is exactly Hypothesis~\ref{h:H2}.
\end{remark}

%==========================================================================
\section{The DeTurck Flow Mollifier and the Besov Regularity Obstruction}\label{sec:besov}
%==========================================================================

\subsection{Failure of the Yang--Mills heat flow}

The Yang--Mills flow is $\partial_tA_\nu=-(d_A^*F_A)_\nu$. Expanding,
\begin{align}
(d_A^*F_A)_\nu&=-\sum_\mu\bigl(\partial_\mu F_{\mu\nu}+[A_\mu,F_{\mu\nu}]\bigr)\notag\\
&=-\sum_\mu\bigl(\partial_\mu\partial_\mu A_\nu-\partial_\mu\partial_\nu A_\mu\bigr)+\mathcal N_\nu(A),
\end{align}
where $\mathcal N_\nu(A)=-\sum_\mu\bigl(\partial_\mu[A_\mu,A_\nu]+[A_\mu,\partial_\mu A_\nu-\partial_\nu A_\mu]+[A_\mu,[A_\mu,A_\nu]]\bigr)$ contains at most one derivative. Hence, with $\Delta=\sum_\mu\partial_\mu^2$,
\begin{equation}
\partial_tA_\nu=\Delta A_\nu-\partial_\nu(\partial\!\cdot\! A)-\mathcal N_\nu(A).
\end{equation}
The linearised operator has principal symbol (replace $\partial_\mu\mapsto i\xi_\mu$)
\begin{equation}\label{eq:symbolYM}
\sigma_{\rm YM}(\xi)_{\nu\rho}=-|\xi|^2\delta_{\nu\rho}+\xi_\nu\xi_\rho=-|\xi|^2\Bigl(\delta_{\nu\rho}-\frac{\xi_\nu\xi_\rho}{|\xi|^2}\Bigr)=-|\xi|^2(\Pi^{\perp}_\xi)_{\nu\rho}.
\end{equation}
Taking $b=\xi$ (the longitudinal direction, i.e.\ the symbol of an infinitesimal gauge orbit $b=d\varphi$, $\hat b=i\xi\hat\varphi$), $\sigma_{\rm YM}(\xi)\xi=-|\xi|^2\xi+\xi|\xi|^2=0$. Thus $\sigma_{\rm YM}(\xi)$ has a kernel of dimension $1$ (times $\dim G$) for every $\xi\neq0$: the flow is only \emph{weakly} parabolic, the heat-kernel smoothing is absent along pure-gauge orbits, and the standard fixed-point theory for the stochastic flow fails.

\subsection{The Yang--Mills--DeTurck flow}

Add the DeTurck term:
\begin{equation}\label{eq:dtflow}
\partial_tB=-d_B^*F_B-d_Bd^*B,\qquad\text{componentwise}\quad \partial_tB_\nu=-(d_B^*F_B)_\nu+D_\nu\psi,\ \ \psi=-d^*B=\sum_\mu\partial_\mu B_\mu.
\end{equation}
(This is exactly the formalised \texttt{IsDeTurckFlow}: the right-hand side is \texttt{-dStarF B + covD B $\mu$ (psi B)}.) At the linear level $-d_Bd^*B\simeq+\partial_\nu\partial_\mu B_\mu$, so
\begin{equation}
\partial_tB_\nu=\Delta B_\nu-\partial_\nu(\partial\!\cdot\!B)+\partial_\nu(\partial\!\cdot\!B)+\text{(lower order)}=\Delta B_\nu+\text{(lower order)},
\end{equation}
and the principal symbol becomes
\begin{equation}
\sigma_{\rm DT}(\xi)=-|\xi|^2\,\mathbf 1_{4}\otimes\mathbf 1_{\g},
\end{equation}
which is strictly negative-definite uniformly on $|\xi|=1$: the flow is \emph{strongly} parabolic. The stochastic version with a $\g$-valued space-time white noise $\xi_{\rm noise}$ is
\begin{equation}\label{eq:dtflow-stoch}
\partial_tB=\Delta B+\mathcal N_{\rm DT}(B)+\xi_{\rm noise},\qquad \E[\xi_{\rm noise}^a(s,x)\xi_{\rm noise}^b(s',y)]=\delta^{ab}\delta(s-s')\delta(x-y),
\end{equation}
where $\mathcal N_{\rm DT}(B)$ collects the terms of the form $B\partial B$ and $B^3$.

\begin{lemma}[Cancellation of the longitudinal drift]\label{lem:drift}
Let $A:\R\to\CalA$ solve the Yang--Mills flow, and let $h:\R\to\CalG$ solve the drift ODE $\partial_th=h\,\psi(A^h)$, $\psi(B)=-d^*B$ (\texttt{IsDriftGauge}). Set $B_t=A_t^{h_t}$. Then
\begin{enumerate}
\item[(i)] $B$ solves the DeTurck flow \eqref{eq:dtflow};
\item[(ii)] for every closed loop $\gamma$ one has $\Tr\,\mathrm{Hol}_\gamma(B_t)=\Tr\,\mathrm{Hol}_\gamma(A_t)$.
\end{enumerate}
\end{lemma}
\begin{proof}
Write $\psi_t=\psi(B_t)$ and suppress $t$ and $x$. Differentiating $h^{-1}h=1$ gives $\partial_t(h^{-1})=-h^{-1}(\partial_th)h^{-1}=-\psi h^{-1}$. Then
\begin{align*}
\partial_tB_\mu&=\partial_t\bigl(h^{-1}A_\mu h+h^{-1}\partial_\mu h\bigr)\\
&=-\psi h^{-1}A_\mu h+h^{-1}\dot A_\mu h+h^{-1}A_\mu h\psi-\psi h^{-1}\partial_\mu h+h^{-1}\partial_\mu(h\psi)\\
&=h^{-1}\dot A_\mu h+\Bigl[-\psi\,(h^{-1}A_\mu h+h^{-1}\partial_\mu h)+(h^{-1}A_\mu h+h^{-1}\partial_\mu h)\,\psi\Bigr]+\partial_\mu\psi,
\end{align*}
where we used $h^{-1}\partial_\mu(h\psi)=h^{-1}(\partial_\mu h)\psi+\partial_\mu\psi$ and regrouped the two $h^{-1}\partial_\mu h$ terms. The bracket is $[B_\mu,\psi]$, hence
\begin{equation*}
\partial_tB_\mu=h^{-1}\dot A_\mu h+\partial_\mu\psi+[B_\mu,\psi]=h^{-1}\dot A_\mu h+D^{B}_\mu\psi.
\end{equation*}
By gauge covariance $h^{-1}\dot A_\mu h=-h^{-1}(d_A^*F_A)_\mu h=-(d_B^*F_B)_\mu$. This is (i).

(ii) Let $U$ solve $U'(s)=-\bigl(\sum_\mu\dot\gamma^\mu A_\mu(\gamma(s))\bigr)U(s)$, $U(0)=1$, and put $V(s)=h(\gamma(s))^{-1}U(s)h(\gamma(0))$. Then, with $\dot\gamma\cdot\partial h=\frac{d}{ds}h(\gamma(s))$,
\begin{align*}
V'&=-h^{-1}(\dot\gamma\!\cdot\!\partial h)h^{-1}Uh(\gamma(0))-h^{-1}(\dot\gamma\!\cdot\!A)Uh(\gamma(0))\\
&=-\bigl[h^{-1}(\dot\gamma\!\cdot\!A)h+h^{-1}(\dot\gamma\!\cdot\!\partial h)\bigr]V=-(\dot\gamma\!\cdot\!B)\,V,
\end{align*}
and $V(0)=1$; by uniqueness for linear ODEs $V=\mathrm{Hol}_\gamma(B)$. For a closed loop on $\T^4$, $\gamma(1)=\gamma(0)+n$ with $n\in\Z^4$, so periodicity of $h$ gives $h(\gamma(1))=h(\gamma(0))=:h_0$ and
\begin{align*}
\mathrm{Hol}_\gamma(B)&=h_0^{-1}\,\mathrm{Hol}_\gamma(A)\,h_0,\\
\Tr\,\mathrm{Hol}_\gamma(B)&=\Tr\bigl(h_0^{-1}\mathrm{Hol}_\gamma(A)h_0\bigr)=\Tr\bigl(\mathrm{Hol}_\gamma(A)h_0h_0^{-1}\bigr)=\Tr\,\mathrm{Hol}_\gamma(A)\end{align*}
by cyclicity of the trace. Hence the longitudinal drift $\psi=-d^*B$ is invisible to all gauge-invariant Wilson-loop observables.
\end{proof}

\begin{remark}
The two statements of the lemma are what the Lean file records as the proposition \texttt{Lemma22} (gauge invariance of \texttt{WilsonTr} and compatibility of the drift gauge with the DeTurck flow). The file supplied with this project contains the \emph{statement} only; no proof term is included, so the Lean formalisation should not be read as a verification of the above argument.
\end{remark}

\subsection{Paraproduct failure at the critical threshold}\label{sec:parap}

Let $\mathcal C^\alpha=B^\alpha_{\infty,\infty}$ with parabolic scaling in $d+1=5$ dimensions (four space dimensions plus flow time). White noise has regularity $\xi_{\rm noise}\in\mathcal C^{-(d+2)/2-\varepsilon}=\mathcal C^{-3-\varepsilon}$, so the stochastic convolution $\Psi=\int_0^t K_{t-s}*\dd W_s$ (gaining two derivatives) lies in $\mathcal C^{-1-\varepsilon}$. Bony's decomposition of a product of $f\in\mathcal C^\alpha$, $g\in\mathcal C^\beta$ is
\begin{equation}
fg=T_fg+T_gf+\Pi(f,g),\qquad T_fg\in\mathcal C^{\min(\beta,\alpha+\beta)},\ \Pi(f,g)\in\mathcal C^{\alpha+\beta}\ \text{ only if }\alpha+\beta>0.
\end{equation}
The quadratic term in $\mathcal N_{\rm DT}$ of type $\Psi\Psi$ has $\alpha=\beta=-1-\varepsilon$, so
\begin{equation}
\alpha+\beta=-2-2\varepsilon<0:
\end{equation}
the resonant product $\Pi(\Psi,\Psi)$ is not defined classically, and the product of the derivative term $\Psi\,\partial\Psi$ is worse still ($-3-2\varepsilon$). In $d=4$ the equation is \emph{critical} (the coupling is dimensionless), so every term of the expansion sits at the same threshold and no subcritical regularity-structure power counting is available; renormalisation is needed already at the level of the first Picard iterate. This is the origin of Hypothesis~\ref{h:H1}: the existence of the renormalised products $\mathcal X$ with $\varepsilon$-uniform bounds is assumed, not proved.

%==========================================================================
\section{BPHZ Renormalisation and Empirical Lattice Verification}\label{sec:bphz}
%==========================================================================

\subsection{The subtraction constant}

For the heat kernel $K_t(x)=(4\pi t)^{-2}e^{-|x|^2/4t}$ on $\R^4$ and the noise \eqref{eq:dtflow-stoch}, the covariance of the stochastic convolution is
\begin{equation}
\E[\Psi^a_t(x)\Psi^b_t(y)]=\delta^{ab}\int_0^tK_{2r}(x-y)\,\dd r,
\end{equation}
using $\E[(K_{t-s}*\xi)(x)(K_{t-s}*\xi)(y)]=K_{2(t-s)}(x-y)$ and the semigroup property. Correlating $\Psi$ with $\Delta\Psi$ at one time gives
\begin{equation}
\E[\Psi^a_t(x)\,\Delta\Psi^b_t(y)]=\delta^{ab}\int_0^t\Delta K_{2r}(x-y)\,\dd r.
\end{equation}
Since $\partial_rK_{2r}=2\Delta K_{2r}$ (because $\partial_sK_s=\Delta K_s$) and $K_{0}=\delta_0$,
\begin{equation}\label{eq:Ct}
\boxed{\;C_t=\int_0^t\Delta K_{2r}\,\dd r\otimes I_{\g}=\tfrac12\bigl(K_{2t}-\delta_0\bigr)\otimes I_{\g}\;}
\end{equation}
with $I_{\g}$ the identity on $\g$. At coincident points the continuum value is $K_{2t}(0)=(8\pi t)^{-2}=1/(64\pi^2t^2)$, and $\delta_0(0)$ is replaced by $a^{-4}$ on a lattice of spacing $a$. The BPHZ counterterms are obtained by subtracting the divergent part of the expected value of each tree-indexed product evaluated against $C_t$.

\subsection{Counterterm classification and the Ward condition}

Power counting: $[B]=1$, $[\partial]=1$, and the right-hand side of the flow has dimension $3$. Local polynomial counterterms $\mathcal C(B)$ compatible with hypercubic symmetry, global adjoint covariance $\mathcal C(g^{-1}Bg)=g^{-1}\mathcal C(B)g$ for constant $g\in G$, and dimension $\le3$ belong to the span of
\begin{align*}
&\text{dim }1:\ B_\mu;\\
&\text{dim }2:\ [B_\nu,B_\mu]\ \text{(antisymmetric, no hypercubic scalar contraction)};\\
&\text{dim }3:\ \Delta B_\mu,\ \partial_\mu\partial\!\cdot\!B,\ [B_\nu,\partial_\nu B_\mu],\ [B_\nu,\partial_\mu B_\nu],\ [\partial\!\cdot\!B,B_\mu],\ [B_\nu,[B_\nu,B_\mu]].
\end{align*}
The terms of dimension $1$ and $2$ carry coefficients of mass dimension $2$ and $1$ and are power divergent; those of dimension $3$ carry logarithmic coefficients.

\begin{definition}[Ward condition $\mathrm W$]\label{def:ward}
A family of counterterms $\mathcal C_\epsilon$ satisfies $\mathrm W$ if $\mathcal C_\epsilon(B^g)=g^{-1}\mathcal C_\epsilon(B)g$ for all $g\in\CalG$ (local, $x$-dependent), up to the DeTurck compensation, i.e.\ the renormalised flow $\partial_tB=-d_B^*F_B-d_Bd^*B+\mathcal C_\epsilon(B)$ is mapped by the drift gauge of Lemma~\ref{lem:drift} onto the renormalised Yang--Mills flow.
\end{definition}

\begin{proposition}[Mass term vanishes]\label{prop:mass}
Suppose $\mathcal C(B)=c_1B$ satisfies the covariance in $\mathrm W$ for $x$-dependent $g$. Then $c_1=0$.
\end{proposition}
\begin{proof}
Compute $\mathcal C(B^g)=c_1(g^{-1}Bg+g^{-1}dg)$, whereas $g^{-1}\mathcal C(B)g=c_1g^{-1}Bg$. The difference is $c_1g^{-1}dg$, which is nonzero for any non-constant $g$ (e.g.\ $g=\exp(\phi T_a)$ with $\phi$ non-constant). Hence $c_1=0$.
\end{proof}

\begin{proposition}[Admissible counterterms under $\mathrm W$]\label{prop:class}
Under $\mathrm W$ the admissible counterterms of dimension $\le3$ form the span of the two covariant objects $d_B^*F_B$ and, modulo the DeTurck compensation, $d_Bd^*B$; consequently the renormalised flow is $\partial_tB=-(1+c_a)\,d_B^*F_B-(1+c_b)\,d_Bd^*B$ with $c_b=c_a$ forced by the compatibility with Lemma~\ref{lem:drift}.
\end{proposition}
\begin{proof}
The dimension-$1$ and (by the same inhomogeneity argument applied to $[B_\nu,B_\mu]$, whose gauge variation contains $[g^{-1}\partial_\nu g,\,g^{-1}B_\mu g]$ without a matching term) the dimension-$2$ monomials are excluded as in Proposition~\ref{prop:mass}, up to combinations that assemble into $F_{\mu\nu}$-type objects, which have dimension $2$ and appear at dimension $3$ only inside $d_B^*F_B$. Among the dimension-$3$ monomials, requiring the non-covariant parts to cancel fixes the relative coefficients of $\Delta B_\mu$, $\partial_\mu\partial\!\cdot\!B$, $[B_\nu,\partial_\nu B_\mu]$, $[B_\nu,\partial_\mu B_\nu]$, $[B_\nu,[B_\nu,B_\mu]]$ to those occurring in the expansion of $(d_B^*F_B)_\mu$ (Section 2.1), and those of $\partial_\mu\partial\!\cdot\!B$, $[\partial\!\cdot\!B,B_\mu]$ to those of $(d_Bd^*B)_\mu$, leaving the two free coefficients $c_a,c_b$. The relation $c_a=c_b$ is the requirement that the drift of Lemma~\ref{lem:drift} still be $-d^*B$ after renormalisation. Completeness of this reduction uses $\mathrm W$ and is therefore conditional on Hypothesis~\ref{h:H3}.
\end{proof}

\subsection{Lattice Fourier simulation}

We evaluated the lattice analogue of \eqref{eq:Ct} on a $16^4$ periodic grid via the four-dimensional FFT: the discrete heat kernel $\hat K_t(k)=\exp(-t\hat\lambda(k))$ with $\hat\lambda(k)=a^{-2}\sum_\mu4\sin^2(ak_\mu/2)$, $k\in(2\pi/16a)\{0,\dots,15\}^4$, the constant $C_t=\tfrac12\bigl(16^{-4}\sum_k\hat K_{2t}(k)-a^{-4}\bigr)$ taken at the origin, and the magnitude $|C_t|$ recorded over a range of flow times. The reported fit of $|C_t|$ to the model
\begin{equation}\label{eq:fit}
|C_t|\;\approx\;m\,\log\!\bigl(1/\sqrt t\bigr)+b
\end{equation}
yielded
\begin{equation}\label{eq:simdata}
m_{\rm sim}=-0.092783\quad(\text{isolated logarithmic slope}),\qquad b_{\rm sim}=0.350591\quad(\text{theoretical convergence intercept}),
\end{equation}
so that the fitted line crosses zero at $\log(1/\sqrt t)=b_{\rm sim}/|m_{\rm sim}|\approx3.779$. These values are quoted exactly as reported by the simulation run supplied with this project; I have not independently re-run it. Figure~\ref{fig:ct} shows the data and fit.

\begin{figure}[ht]
\centering
\includegraphics[width=0.85\textwidth]{fig1.png}
\caption{Logarithmic UV behaviour of the subtraction constant $C_t$ in $d=4$ on the $16^4$ lattice; dashed line: fit \eqref{eq:fit} with slope $m_{\rm sim}$.}
\label{fig:ct}
\end{figure}

\begin{remark}[Scope and limitations of the numerical evidence]\label{rem:limits}
Three qualifications are necessary and are stated here deliberately.
\begin{enumerate}
\item \emph{Quality of the fit.} Figure~\ref{fig:ct} shows that the data are strongly convex in $\log(1/\sqrt t)$: they decay from $\approx0.4$ to $\approx0$ and remain positive beyond the zero crossing $3.779$ of the line, where the line becomes negative. The straight line therefore describes the data only coarsely, and $m_{\rm sim}$ is a window-dependent average slope, not a locally isolated logarithmic coefficient. A claim of logarithmic divergence requires a window with a stable local slope, a goodness-of-fit statistic, and a restriction to $a^2\ll t\ll L^2$; flow times below the lattice scale are dominated by discretisation artefacts.
\item \emph{Comparison with continuum.} The continuum formula \eqref{eq:Ct} predicts a power divergence $\tfrac12(64\pi^2t^2)^{-1}$ in the coincident-point value, not a logarithm. The logarithmic structure of the framework resides in the coefficients $c_a,c_b$ of the marginal (dimension~$3$) counterterms of Proposition~\ref{prop:class}, and the present measurement does not by itself isolate those.
\item \emph{Numerical coincidence.} For $SU(2)$, $2b_0=22/(24\pi^2)=0.092878\ldots$ (with $b_0=11N/(3(4\pi)^2)$ at $N=2$) differs from $|m_{\rm sim}|$ by about $0.1\%$. Given the quality of the fit in item~1, this agreement should be treated as suggestive at most; it is not evidence of asymptotic freedom in the sense of \eqref{eq:2loop}.
\end{enumerate}
Consequently the data \eqref{eq:simdata} are recorded as a consistency check of the lattice implementation and are not used in any hypothesis or proof.
\end{remark}

%==========================================================================
\section{Non-Triviality Witnesses and Exclusion of Degenerations}\label{sec:nogo}
%==========================================================================

A continuum limit must be shown not to collapse to one of three degenerate classes: (a) scale-free, (b) ultralocal (Haar-factorised), (c) Gaussian/free-field. We give one criterion for (a) and one for (b)--(c).

\subsection{Scale-response witness}

Let $w_\rho(R)$ denote the expectation of the renormalised (perimeter- and corner-subtracted, e.g.\ Creutz-ratio) $R\times R$ Wilson loop in the representation $\rho$, with $w_\rho(R)\neq0$. The dimensionless response function is
\begin{equation}
\Gamma_\rho(R)=-\frac{\dd}{\dd\log R}\,\log|w_\rho(R)|.
\end{equation}

\begin{proposition}[Exclusion of scale-free limits]\label{prop:scalefree}
Suppose along the trajectory \eqref{eq:2loop} that $\Gamma_\rho(R)=a_1(\rho)g^2(R)+a_2(\rho)g^4(R)+O(g^6)$ as $R\downarrow0$, with $a_1(\rho)\neq0$ (for $G=SU(N)$ the leading coefficient is proportional to $C_2(\rho)$). Then $\Gamma_\rho$ is not constant on any interval $(0,R_0)$; in particular the limit is not scale-free.
\end{proposition}
\begin{proof}
A scale-free (conformal or power-law) limit has $|w_\rho(R)|=cR^{-\eta}$ and thus $\Gamma_\rho\equiv\eta$, so $\dd\Gamma_\rho/\dd\log R\equiv0$. Using $\dd g^2/\dd\log R=2b_0g^4+O(g^6)$,
\begin{equation}
\frac{\dd\Gamma_\rho}{\dd\log R}=a_1\,\frac{\dd g^2}{\dd\log R}+O(g^6)=2b_0a_1(\rho)\,g^4(R)+O(g^6)\neq0
\end{equation}
for $R$ small, since $b_0>0$, $a_1\neq0$ and $g^2(R)\sim1/(2b_0\log(1/R\Lambda))>0$. This contradicts constancy.
\end{proof}

\subsection{Cumulant witness}

For loops $\gamma_1,\dots,\gamma_4$ and $w_i=\Tr\,\mathrm{Hol}_{\gamma_i}$ define the fourth joint connected cumulant
\begin{equation}
\kappa_4^c(\gamma_1,\dots,\gamma_4)=\E[w_1w_2w_3w_4]-\cdots,
\end{equation}
explicitly $\kappa_4^c=\E[w_1w_2w_3w_4]-\sum_{i}\E[w_i]\E[\prod_{j\ne i}w_j]-\sum_{\{ij\}\{kl\}}\E[w_iw_j]\E[w_kw_l]+2\sum_{\{ij\}}\E[w_iw_j]\E[w_k]\E[w_l]-6\prod_i\E[w_i]$ (sums over the evident index sets, with the standard moment--cumulant coefficients).

\begin{definition}[Criterion $\mathrm K$]
Loops of diameter at most $\ell$ whose mutual distances are $\ge r\ge r_0$ satisfy
\begin{equation}\label{eq:K}
c_-\,e^{-m_+r}\ \le\ |\kappa_4^c(\gamma_1,\dots,\gamma_4)|\ \le\ c_+\,e^{-m_-r},\qquad 0<c_-\le c_+,\ 0<m_-\le m_+ .
\end{equation}
\end{definition}

\begin{remark}
The upper bound is a clustering statement of mass-gap type and the lower bound asserts persistent non-Gaussian correlation. Both are \emph{witness criteria to be verified} for the limit measure; this paper does not prove \eqref{eq:K}. What is proved below is that the three degenerate classes violate it.
\end{remark}

\begin{proposition}[Exclusion of ultralocal Haar factorisation]
Let the limit be ultralocal: links on distinct edges are independent with Haar law, so that Wilson loops with disjoint link supports are independent. Then $\kappa_4^c=0$ for any four loops with pairwise disjoint supports, contradicting the lower bound in \eqref{eq:K}.
\end{proposition}
\begin{proof}
Joint cumulants of independent families vanish, since the cumulant generating function of a sum of independent blocks is additive and $\kappa_4^c$ is the mixed fourth derivative at $0$, which vanishes on any generating function of the form $\sum_i\phi_i(s_i)$. For $r\ge r_0>\ell$ the supports are disjoint, so $\kappa_4^c\equiv0<c_-e^{-m_+r}$.
\end{proof}

\begin{proposition}[Exclusion of power-law Gaussian free-field classes]
Let the loop observables be (nonlinear) functionals of a Gaussian field $\mathcal A$ whose covariance decays like $|x-y|^{-p}$, $p>0$, as for a massless free field. Then either the connected cumulants contributing to the criterion vanish, contradicting the lower bound, or they decay as a power of $r$, contradicting the upper bound in \eqref{eq:K}.
\end{proposition}
\begin{proof}
If $w_i=\langle\mathcal A,f_i\rangle$ are linear in $\mathcal A$, then all joint cumulants of order $\ge3$ vanish by the Gaussian property, contradicting $|\kappa_4^c|\ge c_-e^{-m_+r}>0$. If $w_i=\Phi(\langle\mathcal A,f_i\rangle)$ is a nonlinear local function, expand in Wick monomials. The leading contribution to $\kappa_4^c$ that links all four loops is a tree/loop graph containing at least three propagators between distinct loops; each propagator is bounded below by $c\,r^{-p}$ up to constants in a nondegenerate configuration, so $|\kappa_4^c|\ge c'\,r^{-3p}$ for $r$ large. But $r^{-3p}\le c_+e^{-m_-r}$ fails for large $r$, contradicting the upper bound.
\end{proof}

\begin{theorem}[Conditional non-triviality]
Assume \textup{(H1)--(H3)}, so that a limit measure $\mu$ on $W_L$ exists as a subsequential limit of the renormalised DeTurck lattice measures, and assume the leading coefficient $a_1(\rho)\neq0$ and criterion $\mathrm K$. Then $\mu$ is neither scale-free, nor ultralocal, nor in the power-law Gaussian free-field class.
\end{theorem}
\begin{proof}
Immediate from Propositions~\ref{prop:scalefree} and the two propositions above.
\end{proof}

\subsection{Attribution of the witnesses}
The scale-response function $\Gamma_\rho$ was proposed in the research pipeline by the ChatGPT assistant and the two-sided cumulant bounds by the Perplexity assistant; the derivations in this section were written out and are unreviewed. Neither assistant is an author, and all content requires independent verification by the human authors.

\end{document}
